
\Problem{Nombres semiprimers}

\Statement

En matemàtiques, un natural és \emph{semiprimer} si és el
producte de dos nombres primers. Els primers nombres semiprimers són
4($= 2\times2$), 6($= 3\times2$), 9($= 3\times3$), 10($= 2\times5$), 14($= 2\times7$).

\medskip

1.~ Feu una funció

\begin{lstlisting}
    bool primer(int n);
\end{lstlisting}

que indiqui si un natural $n$ és o
no primer. En aquest apartat, no podeu utilitzar bucles, per tant, serà bo que definiu una
funció auxiliar amb algun paràmetre addicional.

\medskip

2.~ Utilitzeu la funció anterior per implementar una acció

\begin{lstlisting}
    void semiprimer(int n, bool& b, int& p, int& q);
\end{lstlisting}

que, donat un natural
$n$, desa en $b$ si $n$ és o no semiprimer. A més, quan $b$ és cert, desa en
$p$ i $q$ dos nombres primers tals que $pq=n$ i $p \le q$.


\medskip

3.~ Feu un programa que llegeixi una seqüència
de naturals i que, per a cadascun, digui si és o no semiprimer i, en cas afirmatiu,
escrigui la seva descomposició obtinguda per l'acció @semiprimer()@. Seguiu
el format dels exemples.


\Observation

Com sempre en informàtica, els nombres naturals inclouen el zero.

\Sample
